\documentclass[10pt,a4paper]{amsart}
\usepackage[margin=19mm]{geometry}
\usepackage{amsmath,amssymb,mathtools}
\usepackage[scr=rsfs,cal=euler]{mathalfa}
\usepackage{xfrac}
\usepackage[unicode,hidelinks,bookmarks=false]{hyperref}
\newcommand{\ie}{i.e.\ }
\newcommand{\cf}{cf.\ }
\newcommand{\resp}{respectively\ }
\newcommand{\Subsection}{\subsection}
\providecommand{\nobreakdash}{\nobreak\hbox{-}}
\setlength{\emergencystretch}{2em}
\multlinegap0pt
\usepackage{amssymb}
\usepackage[shortlabels]{enumitem}
\usepackage{mathtools}
\newtheorem{definition}{Definition}
\newtheorem{proposition}{Proposition}
\newtheorem{corollary}{Corollary}
\newtheorem{lemma}{Lemma}
\def\A{\mathcal{A}}
\def\C{{\mathbb C}}
\def\Dom{{\mathrm {Dom}}}
\def\L{\mathcal{L}}
\def\N{{\mathbb N}}
\def\d{\bar{\delta}}
\def\db{\bar{\delta}}
\def\pb{\bar{\partial}}
\title{Non-standard Holomorphic Structures on Line Bundles over the Quantum Projective Line}
\author{Mary Graveman, Landen La Rue, Lillian MacArthur, Hunter Pesin, Zhaoting Wei}
\date{Source: arXiv:2510.14263v2 (CC BY 4.0)}
\begin{document}
\maketitle

\noindent\emph{Source and licence.} Non-standard Holomorphic Structures on Line Bundles over the Quantum Projective Line. Version arXiv:2510.14263v2 (CC BY 4.0), distributed by arXiv under the Creative Commons Attribution 4.0 International licence, http://creativecommons.org/licenses/by/4.0/, as recorded in the arXiv licence metadata for this exact version and shown on the abstract page https://arxiv.org/abs/2510.14263v2. Full licence notice: \texttt{licenses/arxiv-2510.14263-CC-BY-4.0.txt}.\par\smallskip

\noindent\textit{Editorial scope.} Counted unit: the article's proposition connections (source0.tex lines 1072--1077). The article's complete original proof of it is delivered with it (source0.tex lines 1078--1132, one \texttt{proof} environment of 55 lines). No mathematics is written, repaired or re-derived: the statement and the proof are the article's own lines, in the article's own notation, and no retained line is substituted or re-flowed.  Retained uncounted as the source-local context the proof needs: lines 618--632; lines 715--717; lines 842--853; lines 884--886; lines 916--924; lines 963--967; lines 1043--1049.  Nothing was omitted from any retained span, so the package carries no omission record.  No retained line contains a citation, so the wrapper carries no bibliography and no bibliography entry.  No retained line contains a figure environment, an image-inclusion command or drawing code, so no image is delivered and the package declares no asset dependency.  Editorial formatting only, in addition: an AMS \texttt{amsart} preamble that restates the article's own declarations and macros used by the retained lines, namely the environments \texttt{definition}, \texttt{proposition}, \texttt{corollary}, \texttt{lemma} on separate counters, together with the standard packages those lines need.  The retained environments print in this document as Definition 1, Proposition 1, Corollary 1, Lemma 1 in order of appearance, whereas the article prints the same items with the numbering of its own full document.  The complete native source of the article as distributed in the arXiv e-print for this exact version is bundled unchanged as \texttt{sources/187171/source0.tex}.
\begin{definition}\label{def: gauge equivalence}
     We call two $\pb$-connections $\overline{\nabla}_{\theta_1}, \overline{\nabla}_{\theta_2}$ on  $\L_n$   gauge equivalent if there exists an invertible element $g\in \mathcal{C}(\C P^1_q)^{\times}\cap \Dom(\pb)$ such that $g^{-1}\in \Dom(\pb)$ and
     $$
         \theta_1=g\theta_2 g^{-1}-\pb(g)g^{-1}.
    $$
     hence
      \begin{equation}\label{eq: gauge equivalence general}
        \pb(g)= g\theta_2 -\theta_1 g.
     \end{equation}

     In particular, a $\pb$-connection $\overline{\nabla}_{\theta}$ on  $\L_n$ is gauge equivalent to the standard $\pb$-connection if there exists an invertible element $g\in \mathcal{C}(\C P^1_q)^{\times}\cap \Dom(\pb)$ such that $g^{-1}\in \Dom(\pb)$ and
     \begin{equation}\label{eq: gauge equivalence to standard general}
        \pb(g)=g\theta.
     \end{equation}
\end{definition}

\begin{equation}\label{eq: mc}
    m_c(f)(x):=f(cx),
\end{equation}

\begin{proposition} \label{prop: delta_bar_domain}
    $f\in C(sp(B_0))$ is in the domain of $\d$ if and only if \[
    \frac{f(x)-f(q^2 x)}{x-q^2 x}
    \]
    is a continuous function on $sp(B_0),$ i.e.
    $$
    \lim_{k\to \infty}\frac{f(q^{2k})-f(q^{2k+2})}{q^{2k}-q^{2k+2}} \text{ exists.}
    $$
In this case we have  \begin{equation}
    (\db f)(q^{2k})=\frac{f(q^{2k})-f(q^{2k+2})}{q^{2k}-q^{2k+2}} \text{ and }(\db f)(0)= \lim_{k\to \infty}\frac{f(q^{2k})-f(q^{2k+2})}{q^{2k}-q^{2k+2}}.
\end{equation}
\end{proposition}

\begin{corollary}\label{coro: inverse domain d}
    If  $f\in C(sp(B_0))$ is in the domain of $\d$ and $f$ is invertible, then $f^{-1}$ is also in the domain of $\d$.
\end{corollary}

\begin{proposition} \label{prop: solution g in product form}
    Let $f\in C(sp(B_0))$ be a function contained in $\Dom(\db).$ Then, a solution to $\d g = g \d f$ is given by \begin{equation}\label{eq: def of g}
    g(q^{2n}) = g(1) \prod_{k = 1}^n  (1 - f(q^{2k - 2}) + f(q^{2k})) \text{ for any }n\geq 1,
    \end{equation}
    and \begin{equation}\label{eq: def of g0}
    g(0) = g(1)\prod_{k = 1}^\infty (1 - f(q^{2k - 2}) + f(q^{2k}))
    \end{equation}
   The $g$ defined above is always in $\Dom(\db)$.
\end{proposition}

\begin{corollary}\label{coro: g invertible}
   The solution $g$ in Proposition \ref{prop: solution g in product form} is invertible if and only if $g(1)\neq 0$ and
    \begin{equation}\label{eq: defective spot of f}
    f(q^{2n}) - f(q^{2n - 2}) \neq 1\text{ for all }n \in \N.\end{equation}
\end{corollary}

\begin{lemma} \label{lemma: commuting partials in CB0}
    For $g, h \in \Dom(\pb)\cap C^*(1, B_0)$ we have  
    \begin{equation}\label{eq: commuting partials in CB0}
    \pb h \cdot g = (m_{q^2}g) \cdot \pb h,
    \end{equation}
    where $m_{q^2}$ is the dilation map in \eqref{eq: mc} extended to $C^*(1, B_0)$ via the functional calculus isomorphism.
\end{lemma}

\begin{proposition}\label{prop: gauge equivalent between connections}
    Let $f, h \in \A(\C P^1_q)\cap C^*(1, B_0),$ and write $S_f, S_h \subset \N$ for the sets of defective spots of $f, h$ respectively. Then, there exists an invertible $g \in C^*(1, B_0)^{\times}\cap \Dom(\db)$ such that \begin{equation}\label{eq: gauge equivalence in B0}
    \pb g = g \pb f -\pb h \cdot g\end{equation} if and only if $S_f = S_h.$

    In particular if $S_f = S_h$, then the two $\pb$-connections $\overline{\nabla}_{\pb f}$ and $\overline{\nabla}_{\pb h}$  are gauge equivalent.
\end{proposition}

\begin{proof}
The second assertion follows from the first one together with  Definition \ref{def: gauge equivalence} and Corollary \ref{coro: inverse domain d}.

    To prove the first assertion, we again use the functional calculus isomorphism $\Psi$ to identify $C^*(1, B_0)$ and $C(sp(B_0))$. Equation \eqref{eq: gauge equivalence in B0} then becomes $$\d g = g \d f - (\d h) g.$$
    By Lemma \ref{lemma: commuting partials in CB0} it becomes 
    \begin{equation}\label{eq: gauge equivalent dilated}
    \d g = g \d f - (m_{q^2} g) \d h,
    \end{equation}
    By Proposition \ref{prop: delta_bar_domain}  we can write \eqref{eq: gauge equivalent dilated}  as \begin{equation}
    \frac{g(x) - g(q^2x)}{x - q^2x} = g(x) \frac{f(x) - f(q^2x)}{x - q^2x} - g(q^2x) \frac{h(x) - h(q^2x)}{x - q^2x}
    \end{equation}
    Since $x - q^2x$ is never zero for $x \in (0,1],$ this becomes \[
    g(x) - g(q^2x) = g(x) (f(x) - f(q^2x)) - g(q^2x) (h(x) - h(q^2x))
    \]
    hence \begin{equation}\label{eq: recursive eq for g}
    g(q^2x)[1 - h(x) + h(q^2x)] = g(x)[1 - f(x) + f(q^2x)]
    \end{equation}
    For $x=q^{2n}$, \eqref{eq: recursive eq for g} becomes
    \begin{equation}\label{eq: recursive eq for g 2}
    g(q^{2n+2})[1 - h(q^{2n}) + h(q^{2n+2})] = g(q^{2n})[1 - f(q^{2n}) + f(q^{2n+2})]
    \end{equation}

If $S_f\neq S_h$, then there must exist an $n$ such that one of $1 - f(q^{2n}) + f(q^{2n+2})$ and $1 - h(q^{2n}) + h(q^{2n+2})$ is zero and the other is non-zero, hence one of $g(q^{2n})$ and $g(q^{2n+2})$ must be zero. Therefore $g$ cannot be invertible.

    If $S_f= S_h$, then $1 - f(q^{2n}) + f(q^{2n+2})$ and $1 - h(q^{2n}) + h(q^{2n+2})$ are both zero or both nonzero. If both are nonzero, then we have 
    \begin{equation}
    g(q^{2n+2})=g(q^{2n})\frac{1 - f(q^{2n}) + f(q^{2n+2})}{1 - h(q^{2n}) + h(q^{2n+2})}.
    \end{equation}
    If both are zero, then \eqref{eq: recursive eq for g 2} implies that $g(q^{2n+2})$ can be any number. We can therefore define $g$ inductively at any $q^{2n}$ so that $g(q^{2n})\neq 0$.

    

   Moreover since $S_f=S_h$ is a finite set, let 
   $$
   N= \text{the maximum of }S_f.
   $$
Then for any $n>N$, $1 - f(q^{2n}) + f(q^{2n+2})$ and $1 - h(q^{2n}) + h(q^{2n+2})$ are both nonzero hence $g(q^{2n})$ is uniquely determined by $g(q^{2N})$ by
\begin{equation}
    g(q^{2n}) = g(q^{2N})\prod_{k=N}^{n-1} \frac{1 - f(q^{2k}) + f(q^{2k+2})}{1 - h(q^{2k}) + h(q^{2k+2})} .
    \end{equation}

    Therefore
  \begin{equation}\label{eq: infinity product quotient}
    \lim_{n \to \infty} g(q^{2n}) 
  =\prod_{k = N }^{\infty}  \frac{1 - f(q^{2k}) + f(q^{2k + 2})}{1 - h(q^{2k}) + h(q^{2k + 2})}
    \end{equation}

   Since $f,h\in \Dom(\db)$, by the same argument as in the proof of Proposition \ref{prop: solution g in product form} and Corollary \ref{coro: g invertible}, the infinite product on the right hand side of \eqref{eq: infinity product quotient} converges with a nonzero limit. Thus, $g(0)$ exists, and is nonzero.
   Hence $g$ is a continuous function which is invertible.
    
    It remains to show that $g\in \Dom(\db)$. But this  follows from \[
    \d g = g \d f - \d h \cdot g,
    \]
    and the fact that $\d f, \d h,$ and $g$ are all continuous.
\end{proof}

\end{document}
